\section{Soft truncation with an approximate energy origin}

For a Hermitian matrix $A$, write $T_s(A)=\min(A,sI)$ for its hard
truncation at the energy $s$, defined by functional calculus.

\begin{theorem}[Monotonicity in the hard cutoff]
    \label{thm:hard_truncation_cutoff_monotone}
    \lean{Matrix.monotone_energyTruncation}
    \leanok
    If $s\leq t$, then $T_s(A)\leq T_t(A)$.
\end{theorem}
\begin{proof}\leanok
    \uses{def:energy_truncation}
    The scalar inequality $\min\{x,s\}\leq\min\{x,t\}$ holds at every
    spectral value of $A$. Apply order preservation of the functional
    calculus for this fixed matrix.
\end{proof}

\begin{theorem}[Translation of the hard cutoff]
    \label{thm:hard_truncation_energy_shift}
    \lean{Matrix.IsHermitian.energyTruncation_sub_smul_one}
    \leanok
    For all real $r,t$, one has
    \begin{align}
        T_t(A-rI)=T_{r+t}(A)-rI.
        \label{eq:hard_truncation_energy_shift}
    \end{align}
\end{theorem}
\begin{proof}\leanok
    \uses{def:energy_truncation}
    Composition of the functional calculus reduces the assertion to
    $\min\{x-r,t\}=\min\{x,r+t\}-r$.
\end{proof}

\begin{theorem}[Soft truncation with an approximate energy origin]
    \label{thm:soft_hard_truncation_energy_shift}
    \lean{Matrix.IsHermitian.softTruncation_shifted_hard_bounds}
    \leanok
    Let $r,r'$ be real numbers such that $A-r'I\geq0$ and $r-10\leq r'$.
    For $t>0$ and an integer $m\geq1$, put $\varepsilon=(t/2)2^{-m}$.
    Then
    \begin{align}
        T_{r+t/2-10}(A)-\varepsilon I
        \leq r'I+h_{m,t}(A-r'I)\leq A.
        \label{eq:soft_hard_truncation_energy_shift}
    \end{align}
\end{theorem}
\begin{proof}\leanok
    \uses{thm:matrix_soft_hard_truncation,
        thm:hard_truncation_energy_shift,
        thm:hard_truncation_cutoff_monotone}
    Apply~(\ref{eq:matrix_soft_hard_truncation}) to $A-r'I$ and add $r'I$.
    By~(\ref{eq:hard_truncation_energy_shift}), the hard comparator is
    $T_{r'+t/2}(A)$. Since $r+t/2-10\leq r'+t/2$, monotonicity of the
    hard cutoff gives~(\ref{eq:soft_hard_truncation_energy_shift}).
\end{proof}

In the regional soft truncation of
Arad--Landau--Vazirani--Vidick~\cite[lines 1178--1185]{Arad2016RigorousRG},
$r$ is the regional ground energy and $r-10\leq r'\leq r$. Thus
$A-r'I\geq0$. Taking $t=m$ gives the hard threshold $m/2-10$ above the true
regional ground energy. The crossing interactions belong to the unchanged
complementary Hamiltonian; they are not subtracted from this threshold.

\begin{theorem}[Coercivity after one centered soft replacement]
    \label{thm:soft_truncation_centered_coercivity}
    \lean{Matrix.softTruncation_centered_coercivity}
    \leanok
    Let $(A_i)$ be a finite family of Hermitian matrices with $A_i\ge e_iI$,
    and fix an index $j$. Suppose $e_j-10\le r'\le e_j$, $t>0$, and $m\ge1$.
    Put
    \begin{align*}
        F&=r'I+h_{m,t}(A_j-r'I),&
        \alpha&=(t/2)2^{-m},& \tau&=t/2-10,\\
        X&=\sum_{i\ne j}(A_i-e_iI)+(F-e_jI).
    \end{align*}
    For $0\le u\le\tau$, one has
    \begin{align}
        X+\alpha I&\ge0,&
        u\,\mathbf1_{[e_j+u,\infty)}(A_j)&\le X+\alpha I.
        \label{eq:soft_truncation_centered_coercivity}
    \end{align}
    Neither commutativity of the family nor attainment of the lower bounds
    is required for these order inequalities.
\end{theorem}
\begin{proof}\leanok
    \uses{thm:soft_hard_truncation_energy_shift,
        thm:partial_clipping_high_projection}
    The shifted comparison gives $T_{e_j+\tau}(A_j)-\alpha I\le F$.
    Adding the unchanged centered terms shows that the actual centered hard
    sum is at most $X+\alpha I$. The hard high-projection bound applies at
    $u$ and at zero. These give both assertions in
    (\ref{eq:soft_truncation_centered_coercivity}).
\end{proof}

\begin{theorem}[A cutoff satisfying the reverse-weight inequalities]
    \label{thm:soft_truncation_reverse_cutoff}
    \lean{Real.softTruncation_reverse_cutoff_bounds}
    \leanok
    For an integer $m\ge0$, put $\alpha=(m/2)2^{-m}$. Let
    $-\alpha\le E\le2$, $\eta\ge0$, and $0\le C\le4$. Suppose
    \begin{align}
        m/2\ge17+\eta+\alpha,
        \label{eq:soft_truncation_reverse_scale}
    \end{align}
    and set $u=E+\alpha+\eta+C+1$. Then
    \begin{align*}
        0\le u&\le m/2-10,& u&\le7+\eta+\alpha,\\
        C&<u-(E+\alpha+\eta),&
        u-(E+\alpha+\eta)-C&=1.
    \end{align*}
\end{theorem}
\begin{proof}\leanok
    The lower bounds on $E+\alpha$, $\eta$ and $C$ give $u\ge0$.
    The upper bounds on $E$ and $C$ give $u\le7+\eta+\alpha$;
    (\ref{eq:soft_truncation_reverse_scale}) then gives $u\le m/2-10$.
    Cancellation gives the last identity and hence the strict inequality.
\end{proof}

\begin{lemma}[Exponential weights from their high-energy restriction]
    \label{lem:exponential_weight_high_restriction}
    \lean{Matrix.norm_exp_domRestrict_le_of_high_restriction}
    \leanok
    Let $A$ be Hermitian, let $s\ge0$, and put
    $Q=\mathbf1_{[r+u,\infty)}(A)$ and $W=e^{s(A-(r+u)I)}$.
    Let $G$ be a subspace. If $WQ|_G$, viewed as a map into the range of
    $Q$, has norm at most $C$, then
    \begin{align*}
        \|W|_G\|\le1+C.
    \end{align*}
\end{lemma}
\begin{proof}\leanok
    On the complementary spectral space, $W(I-Q)$ has norm at most one.
    Decompose $W=WQ+W(I-Q)$ and use the triangle inequality.
\end{proof}

\begin{lemma}[Changing the origin of an exponential weight]
    \label{lem:shifted_exponential_weight_restriction}
    \lean{Matrix.norm_shifted_exp_domRestrict_le}
    \leanok
    For a finite complex matrix $A$, real numbers $r,u,\varepsilon,s$,
    and a subspace $G$, a bound
    $\|e^{s(A-(r+u)I)}|_G\|\le M$ implies
    \begin{align*}
        \|e^{s(A-\varepsilon I)}|_G\|
          \le e^{s(r+u-\varepsilon)}M.
    \end{align*}
\end{lemma}
\begin{proof}\leanok
    The scalar matrix commutes with $A$, so the two exponentials differ by
    the positive scalar factor $e^{s(r+u-\varepsilon)}$.
\end{proof}

\begin{theorem}[A soft remainder restricted to an arbitrary subspace]
    \label{thm:soft_truncation_exponential_remainder_restriction}
    \lean{Matrix.norm_softTruncation_remainder_domRestrict_le_exp_weight}
    \leanok
    Let $A\ge0$ be a finite complex matrix, let $m\ge1$ be an integer,
    and let $G$ be any subspace. Set $d_m(x)=x-h_{m,m}(x)$. Then
    \begin{align*}
        \|d_m(A)|_G\|
          \le16e^{-m/16}\|e^{A/8}|_G\|.
    \end{align*}
    This matrix estimate does not assume or establish a physical
    exponential-moment bound.
\end{theorem}
\begin{proof}\leanok
    \uses{thm:scalar_soft_truncation_exp_weight}
    The scalar exponential remainder estimate gives
    $0\le d_m(x)e^{-x/8}\le16e^{-m/16}$ for $x\ge0$.
    Functional calculus therefore bounds
    $B_m=d_m(A)e^{-A/8}$ by $\|B_m\|\le16e^{-m/16}$.
    Since $d_m(A)=B_me^{A/8}$, restriction and multiplication give the claim.
\end{proof}
